\documentclass[12pt,letterpaper,oneside]{article}
\usepackage[utf8]{inputenc}
\usepackage[T1]{fontenc}
\usepackage[spanish]{babel}
\usepackage{lmodern}
\usepackage[margin=2.5cm]{geometry}
\usepackage{amsmath,amssymb}
\usepackage{booktabs,tabularx,longtable,array}
\usepackage{graphicx,rotating,lscape,float}
\usepackage[authoryear,round]{natbib}
\usepackage[hidelinks]{hyperref}
\usepackage{caption}
\usepackage{xcolor}
\linespread{1.25}
\setlength{\parindent}{1.25cm}
\setlength{\parskip}{0pt}
\captionsetup{font=small,labelfont=bf}
\renewcommand{\arraystretch}{1.15}
\newcommand{\pval}[1]{\textit{p}~#1}
\newcommand{\NA}{\textit{n.a.}}

\title{Confianza interpersonal, confianza institucional y familiaridad con mecanismos de participación ciudadana: un modelo de ecuaciones estructurales con la ECP 2023}
\author{Juan David Roa Castañeda\\Gabriel Felipe Aldana Vargas\\Laura Camila Rodríguez Granados\\[4pt]\small Facultad de Estadística, Universidad Santo Tomás}
\date{Modelación de Ecuaciones Estructurales --- 2026-II}

\begin{document}
\maketitle

\begin{abstract}
\noindent\textbf{Objetivo:} Estimar la asociación entre la confianza interpersonal, la confianza institucional y la familiaridad autorreportada con mecanismos de participación ciudadana en personas encuestadas en Colombia en 2023. \textbf{Datos y método:} Se integraron cinco módulos de la Encuesta de Cultura Política (ECP) 2023 del DANE mediante cuatro claves individuales; la intersección fue de 46.392 personas. Se estimó un modelo de tres factores con indicadores ordinales y binarios mediante WLSMV, faltantes por pares y ponderación por el factor de expansión \texttt{FEX\_P}. Se evaluó el ajuste global y local, la confiabilidad, la validez y la invariancia ordinal por sexo. \textbf{Resultados:} El ajuste robusto fue débil (CFI=0{,}900; TLI=0{,}890; RMSEA=0{,}092; SRMR=0{,}040), aunque los índices escalados fueron más favorables (CFI=0{,}981; RMSEA=0{,}046). Las rutas hacia familiaridad fueron positivas y pequeñas (\(\beta=0{,}094\) y \(\beta=0{,}065\)); el modelo explicó 1,6\% de su varianza. La asociación institucional fue mayor entre hombres en la comparación multigrupo, mientras que la interpersonal no difirió estadísticamente por sexo. \textbf{Conclusión:} Los datos sugieren asociaciones modestas, pero el ajuste de medición y la varianza no ajustada al diseño limitan la inferencia. Los resultados se consideran exploratorios, no causales ni plenamente representativos en términos de incertidumbre.

\medskip\noindent\textbf{Palabras clave:} confianza interpersonal; confianza institucional; participación ciudadana; ecuaciones estructurales; Colombia.
\end{abstract}

\selectlanguage{english}
\begin{abstract}
\noindent\textbf{Objective:} To estimate the association of interpersonal trust and institutional trust with self-reported familiarity with citizen participation mechanisms among people surveyed in Colombia in 2023. \textbf{Data and method:} Five modules of the 2023 Political Culture Survey (ECP) produced by DANE were linked using four individual identifiers, yielding an intersection of 46,392 people. A three-factor model with ordinal and binary indicators was estimated using WLSMV, pairwise available data, and the \texttt{FEX\_P} expansion weight. Global and local fit, reliability, validity, and ordinal measurement invariance by sex were examined. \textbf{Results:} Robust fit was weak (CFI=0{,}900; TLI=0{,}890; RMSEA=0{,}092; SRMR=0{,}040), whereas scaled indices were more favorable (CFI=0{,}981; RMSEA=0{,}046). Both paths to familiarity were positive and small (\(\beta=0{,}094\) and \(\beta=0{,}065\)); the model explained 1.6\% of the outcome variance. The institutional trust association was larger among men in the multigroup comparison, while the interpersonal trust association did not differ statistically by sex. \textbf{Conclusion:} The data suggest modest associations, but measurement fit and variance estimates that do not account for the full survey design limit inference. Findings are exploratory, noncausal, and not fully design-adjusted for uncertainty.

\medskip\noindent\textbf{Keywords:} interpersonal trust; institutional trust; citizen participation; structural equation modeling; Colombia.
\end{abstract}
\selectlanguage{spanish}

\section{Introducción}

La participación ciudadana requiere que las personas reconozcan canales mediante los cuales pueden intervenir en asuntos públicos, pero el reconocimiento de esos mecanismos no equivale por sí mismo a utilizarlos ni a comprender sus procedimientos. La Encuesta de Cultura Política (ECP) del Departamento Administrativo Nacional de Estadística (DANE) mide, entre otros dominios, la confianza declarada en personas e instituciones y si las personas conocen o han oído hablar de mecanismos de participación. La ronda de 2023 se recolectó entre septiembre y diciembre y sus resultados oficiales se expanden a la población civil no institucional de 18 años o más de los dominios cubiertos por la operación \citep{DANE2024Boletin,DANEECP2023}.

La confianza se ha vinculado a la cooperación cívica, la acción colectiva y las orientaciones frente al sistema político. Sin embargo, la confianza interpersonal y la confianza institucional son conceptos distintos: la primera se refiere a expectativas hacia otras personas, mientras que la segunda expresa una evaluación de actores e instituciones públicas \citep{LeviStoker2000,DelheyNewton2005}. Asimismo, las asociaciones entre confianza y participación pueden ser recíprocas, estar condicionadas por experiencias previas o reflejar selección, lo cual impide dar una interpretación causal a datos transversales \citep{vanIngenBekkers2015,Devine2024}.

Este trabajo analiza si los niveles más altos de confianza interpersonal y confianza institucional se asocian con mayor familiaridad declarada con seis mecanismos formales: revocatoria del mandato, iniciativa popular, referendo, plebiscito, consulta popular y cabildo abierto. El análisis continúa la propuesta aprobada para el Taller 1 y aplica el enfoque de dos pasos que recomienda evaluar primero el modelo de medición y, luego, el modelo estructural \citep{AndersonGerbing1988}. La pregunta es: ¿cómo se asocian la confianza interpersonal y la confianza institucional con la familiaridad autorreportada con mecanismos de participación ciudadana entre personas encuestadas en Colombia en 2023?

El aporte empírico consiste en integrar módulos temáticos de la ECP, incorporar el factor oficial de expansión y evaluar diferencias entre hombres y mujeres con un modelo ordinal multigrupo. El estudio también hace explícita una limitación relevante para cualquier lectura poblacional: aunque se identificó el factor de expansión, los microdatos anonimizados disponibles no contienen los identificadores de estratos ni de las unidades de muestreo requeridos para reproducir plenamente la varianza del diseño. Se reporta, por tanto, un análisis ponderado de carácter exploratorio y se separa la magnitud de asociación de la significación estadística.

\section{Marco teórico e hipótesis}

\subsection{Confianza interpersonal y familiaridad cívica}

La confianza interpersonal facilita interacciones en las que circula información, se construyen expectativas de reciprocidad y se coordinan acciones colectivas. Los enfoques de capital social relacionan redes, normas y confianza con la cooperación cívica, aunque la naturaleza de las redes y la forma en que se mide la confianza pueden modificar estas relaciones \citep{Putnam1993}. La literatura distingue la confianza generalizada hacia personas desconocidas de la confianza particularizada en grupos conocidos; una escala que combina desconocidos, vecinos, amistades y personas de otra nacionalidad comprende componentes de ambos tipos y no debe interpretarse como una medida pura de confianza generalizada \citep{DelheyNewton2005}.

La hipótesis de asociación positiva plantea que las personas que declaran mayor confianza en esos grupos pueden tener más oportunidades de conversar sobre asuntos públicos, participar en redes comunitarias o exponerse a información cívica. No supone que la confianza produzca conocimiento: la participación puede aumentar la confianza, y características individuales no modeladas pueden relacionarse con ambas variables. Esta expectativa se contrasta como asociación transversal.

\noindent\textbf{H1.} Una mayor confianza interpersonal se asocia positivamente con la familiaridad autorreportada con mecanismos de participación ciudadana.

\subsection{Confianza institucional y familiaridad con mecanismos}

La confianza política comprende evaluaciones de la confiabilidad y legitimidad de instituciones y representantes. No es una medida directa de su desempeño ni de la confianza hacia otras personas \citep{LeviStoker2000}. Las personas con mayor confianza institucional podrían prestar más atención a los canales formales de incidencia y reconocer mejor sus reglas; también es plausible que la percepción de instituciones poco confiables se relacione con mayor interés en mecanismos de control o revocatoria. La dirección se define aquí como expectativa positiva por su coherencia con una orientación favorable hacia la vida institucional, pero se reconoce como la hipótesis más incierta.

\noindent\textbf{H2.} Una mayor confianza institucional se asocia positivamente con la familiaridad autorreportada con mecanismos de participación ciudadana.

La confianza interpersonal y la institucional pueden compartir experiencias de cooperación, información y evaluación cívica, pero se mantienen como latentes diferenciadas. Se permite su covarianza libre en el modelo. La asociación empírica entre ambas se reporta como parte de la estructura de covarianzas, sin tratarla como una causa de la familiaridad.

\subsection{Familiaridad declarada y conocimiento político}

Los seis reactivos de la ECP preguntan si la persona conoce o ha oído hablar de distintos mecanismos. Capturan reconocimiento declarado y no una prueba de conocimiento factual, comprensión jurídica, uso efectivo o participación. La literatura de conocimiento político subraya la diferencia entre recordar información factual y otras formas de involucramiento o competencia cívica \citep{DelliCarpiniKeeter1996}. Por precisión conceptual, el factor endógeno se denomina en adelante \emph{familiaridad autorreportada con mecanismos}, aunque la etiqueta interna del modelo y el cuestionario se refieran a conocimiento.

La cultura cívica ubica las orientaciones individuales dentro del sistema político y destaca que la participación depende de actitudes y competencias que se forman en contextos sociales \citep{AlmondVerba1963}. Esta perspectiva motiva analizar conjuntamente confianza y familiaridad, pero no implica que un único modelo reflectivo agote la complejidad de los constructos. La estructura factorial se somete a evaluación empírica antes de interpretar las rutas.

\section{Datos}

\subsection{Fuente y población}

La ECP 2023 es una encuesta probabilística, estratificada, multietápica y de conglomerados. Su población objetivo comprende población civil no institucional de 18 años o más en los territorios cubiertos; la operación excluye varios departamentos de la Amazonía y Orinoquía y el archipiélago de San Andrés, Providencia y Santa Catalina \citep{DANE2024Boletin,DANEECP2023}. El DANE informó 64.770 personas encuestadas en la operación completa \citep{DANE2024Boletin}. La base usada aquí es una intersección temática menor y no se presenta como la muestra completa de la encuesta.

Se integraron los módulos Democracia, Participación, Capital Social, Características generales y Características generales 2 mediante las claves \texttt{DIRECTORIO}, \texttt{NRO\_ENCUESTA}, \texttt{HOGAR\_NUMERO} y \texttt{PERSONA\_NUMERO}. Las claves no tuvieron duplicados ni valores incompletos en los módulos utilizados. La intersección común fue de 46.392 personas. En la base de Características generales se encontraron 64.759 registros, pero solo 46.392 coincidieron con los módulos temáticos requeridos. Esta diferencia de cobertura se debe reportar como condición de la muestra analítica, no como una discrepancia que pueda corregirse con ponderación.

En el Taller 1 se reportó que no había ponderadores porque los cuatro ZIP usados entonces no incluían \texttt{FEX\_P}. En el Taller 2 se enlazó ese factor desde la Tabla de viviendas mediante \texttt{DIRECTORIO}. Todos los registros analíticos recibieron un peso positivo. La suma de pesos fue 37.076.519, cifra que resulta de los registros integrados y no debe confundirse con un total poblacional oficial sin cotejar dominios, elegibilidad y universo. La cobertura y el diseño de la ECP son definidos por el DANE y delimitan el alcance de la interpretación.

\subsection{Indicadores y procesamiento}

La Tabla~\ref{tab:indicadores} resume los constructos e indicadores. Las escalas de confianza institucional e interpersonal tienen cinco categorías ordenadas, donde valores mayores indican más confianza. El código 99 (no sabe/no informa) se convirtió en faltante. Los reactivos P5263S6 (Asamblea departamental) y P5263S12 (Gobernación) tienen no aplicabilidad estructural en Bogotá. Los seis indicadores P5376S1--P5376S6 se recodificaron como 1 si la persona respondió que conoce o ha oído hablar del mecanismo y 0 si respondió que no.

En el perfil inicial, 35.024 personas (75,5\%) respondieron todos los 25 indicadores; las restantes aportan información disponible por pares en la estimación WLSMV. Los seis reactivos de familiaridad no presentaron valores faltantes en la base integrada. En confianza interpersonal, los faltantes por reactivo variaron entre 0,9\% y 2,4\%. En confianza institucional, oscilaron entre 5,4\% y 19,9\%, con las mayores proporciones en Asamblea y Gobernación, debido a la no aplicabilidad y a las respuestas no informadas. No se eliminaron casos con base exclusiva en la regla de distancia chi-cuadrado de normalidad: el perfil previo identificó que esa regla no es adecuada para excluir automáticamente observaciones con indicadores ordinales y binarios.

\begin{table}

\caption{\label{tab:indicadores}Constructos, indicadores y codificación}
\centering
\begin{tabularx}{\textwidth}{p{2.8cm}X p{2.5cm} X}
\toprule
Constructo & Contenido & Indicadores & Tratamiento\\
\midrule
Confianza interpersonal & Confianza en personas desconocidas, vecinos, amistades y personas de otra nacionalidad & P2025S4--P2025S7 & Ordinal 1--5; 99 a faltante\\
Confianza institucional & Confianza en 15 instituciones y actores públicos & P5263S1--P5263S15 & Ordinal 1--5; 99 a faltante; S6/S12 no aplican en Bogotá\\
Familiaridad autorreportada & Reconocimiento de seis mecanismos formales de participación & P5376S1--P5376S6 & Binaria; 1=conoce/ha oído y 0=no conoce\\
\bottomrule
\end{tabularx}
\end{table}

La Tabla~\ref{tab:indicadores} identifica los códigos del formulario. Familiaridad autorreportada no equivale a conocimiento factual ni a uso de los mecanismos. Fuente: DANE, formulario ECP 2023.

\section{Métodos}

\subsection{Especificación e identificación}

Se estimó un modelo de medición confirmatorio con tres factores correlacionados: confianza interpersonal (cuatro indicadores), confianza institucional (quince) y familiaridad con mecanismos (seis). El modelo estructural regresa el factor de familiaridad sobre ambos factores de confianza y deja libre la covarianza de estos predictores:
\begin{align}
  x &= \Lambda_x\xi+\delta, & y &= \Lambda_y\eta+\epsilon, \\
  \eta &= \gamma_1\xi_1+\gamma_2\xi_2+\zeta,
\end{align}
donde \(\xi_1\) es confianza interpersonal, \(\xi_2\) confianza institucional y \(\eta\) familiaridad autorreportada. Se fijó en uno la carga del indicador marcador de cada factor para establecer la escala. La estructura inversa planteada en el Taller 1 es observacionalmente equivalente bajo covarianzas residuales libres entre predictores; por ello el ajuste global no permite elegir una dirección causal \citep{MacCallum1993}.

La cuenta de covarianzas tiene \(t=53\) parámetros libres y 325 momentos, de modo que \(gl=325-53=272\). En la formulación ordinal se estiman 82 umbrales y 382 estadísticas muestrales; con 110 parámetros libres se obtienen igualmente 272 grados de libertad. Esta cuenta coincide con los grados de libertad observados en lavaan y verifica la identificación por grados de libertad, sin constituir por sí sola prueba de identificación local.

\begin{figure}[htbp]
\centering
\begin{center}
\renewcommand{\arraystretch}{1.5}
\begin{tabular}{c c c}
\fbox{\parbox{4cm}{\centering Confianza interpersonal\\4 indicadores; $\lambda_{P2025S4}=0{,}674$}} &
\(\beta=.094^{***}\)\(\longrightarrow\) &
\fbox{\parbox{4cm}{\centering Familiaridad autorreportada\\6 indicadores; $\lambda_{P5376S1}=0{,}857$}}\\
\(r=0{,}265^{***}\) & & \(R^2=0{,}016\)\\
\fbox{\parbox{4cm}{\centering Confianza institucional\\15 indicadores; $\lambda_{P5263S1}=0{,}840$}} &
\(\beta=.065^{***}\)\(\longrightarrow\) &
\fbox{\parbox{4cm}{\centering Familiaridad autorreportada\\6 indicadores; $\lambda_{P5376S1}=0{,}857$}}\\
\end{tabular}
\end{center}
\vspace{-1em}
\begin{center}\footnotesize Nota. ***\(p<.001\). Las rutas representan asociaciones transversales estimadas en el modelo ponderado.\end{center}

\caption{Modelo de medición y rutas estructurales estandarizadas. Las rutas del modelo completo son positivas; la comparación multigrupo posterior evalúa si difieren por sexo. Fuente: estimaciones propias con la ECP 2023 del DANE.}
\label{fig:modelo}
\end{figure}

\subsection{Estimación, faltantes y diseño muestral}

Por la naturaleza ordinal y dicotómica de los indicadores, se empleó WLSMV, estimador que utiliza umbrales y asociaciones policóricas/poliseriales y proporciona estadísticos ajustados para datos categóricos \citep{Li2016,Rosseel2012}. Los faltantes se trataron por pares, conforme a la alternativa WLSMV indicada en el Taller 2. La ausencia estructural de dos ítems institucionales para Bogotá permaneció codificada como faltante. \texttt{FEX\_P} se incorporó como peso de muestreo en los modelos; no se reescaló ni se derivó un factor alternativo.

Aunque la ECP tiene diseño complejo, los ZIP anonimizados públicos usados no contienen identificadores de estrato ni UPM/conglomerado. \texttt{REGION} identifica dominios geográficos y \texttt{DIRECTORIO} una vivienda; no son sustitutos válidos de esas unidades de diseño. El equipo confirmó que el DANE no facilita esos identificadores ni ofrece el cálculo alternativo de varianza para este análisis. En consecuencia, la ponderación informa los puntos estimados, pero los errores estándar, intervalos y pruebas no incorporan plenamente la estratificación y conglomeración. Todas las decisiones inferenciales se califican como exploratorias.

\subsection{Evaluación del ajuste y validez}

Para todos los modelos se revisaron convergencia, soluciones impropias, \(\chi^2\), grados de libertad, CFI, TLI, RMSEA con intervalo de confianza del 90\% y SRMR. Se reportan tanto las variantes robustas como las escaladas de los índices WLSMV, porque difieren de manera sustantiva en estos datos. Los puntos de corte comunes se interpretan como guías y no como reglas mecánicas, dado que el rendimiento de los índices depende del modelo, las categorías y el estimador \citep{HuBentler1999,Marsh2004,XiaYang2019}.

Se examinaron residuos de correlación y los índices de modificación (IM) junto con EPC. No se incorporaron covarianzas residuales ni se eliminaron indicadores solo para mejorar el ajuste; la modificación post hoc capitaliza el azar y requiere fundamento teórico y validación independiente \citep{MacCallum1992,Kline2023}. Se estimó confiabilidad compuesta y varianza media extraída (AVE), junto con el criterio de Fornell--Larcker y HTMT \citep{FornellLarcker1981,Henseler2015}. La función usada para HTMT no incorporó el peso; por ello esos valores se presentan únicamente como diagnóstico descriptivo no ponderado.

\subsection{Hipótesis, magnitud y modelos alternativos}

Las hipótesis de signo positivo se contrastaron unilateralmente y se reportaron B, error estándar, \(z\), \(p\), \(\beta\) estandarizado e intervalo del 95\% para B. Dado el N grande, se priorizó la magnitud y dirección frente a la significación aislada. Se reportaron \(R^2\) del factor endógeno y \(f^2\) incremental de cada predictor. El modelo aprobado no incluye mediador ni segundo factor endógeno, así que los efectos indirectos y totales mediados no aplican a esta especificación.

Se comparó descriptivamente el modelo propuesto con una estructura exploratoria que dividió confianza institucional en dos factores. Los modelos no son anidados y WLSMV no proporcionó AIC/BIC convencionales; no se declaró un ganador exclusivamente por índices. El AFC correlacionado y el SEM de referencia tienen el mismo número de grados de libertad y prácticamente el mismo ajuste: la prueba de diferencia no tiene grados de libertad diferenciales y no permite probar si las relaciones teóricas empeoran el AFC.

\subsection{Invariancia multigrupo}

Se eligió sexo como variable de agrupación, con 21.405 hombres y 24.987 mujeres. Se estimó la secuencia ordinal de identificación Wu--Estabrook, primero configural, luego igualdad de umbrales y después umbrales más cargas, con parametrización theta \citep{WuEstabrook2016,Jorgensen2022}. Se consideraron cambios de CFI, RMSEA y SRMR junto con las pruebas de diferencia ajustadas. Tras la invariancia métrica ordinal, se contrastó la igualdad conjunta de las dos rutas mediante modelos libres e iguales y se evaluó cada restricción con pruebas de Wald. Por la limitación de diseño y el ajuste absoluto, estos contrastes se consideran exploratorios.

\subsection{Software}

El procesamiento se ejecutó con R 4.6.1, \texttt{lavaan} 0.7-3 y \texttt{semTools} 0.5-9 según la sesión reproducida en este equipo \citep{RCore2026,Rosseel2012,Jorgensen2022}. La secuencia de scripts fija una semilla, usa rutas relativas, integra los módulos, ajusta los modelos y genera las tablas que sustentan los resultados.

\section{Resultados}

\subsection{Ajuste global y local del modelo de medición}

El AFC convergió y pasó las comprobaciones de solución admisible: no se observaron varianzas negativas ni correlaciones latentes mayores que uno. El modelo de tres factores obtuvo \(\chi^2_{(272)}=27.181,35\), \(p<{,}001\). Los índices robustos fueron CFI=0{,}900, TLI=0{,}890 y RMSEA=0{,}092, IC 90\% [0{,}091,0{,}093]; el SRMR fue 0{,}040. Los índices escalados fueron CFI=0{,}981, TLI=0{,}979 y RMSEA=0{,}046. La diferencia entre las versiones robustas y escaladas es material; por prudencia, no se interpreta el ajuste como inequívocamente adecuado. El chi-cuadrado se reporta como estadístico ajustado, pero su significación con este tamaño muestral no decide por sí sola la validez.

\begingroup\scriptsize\setlength{\tabcolsep}{2pt}
\begin{table}

\caption{\label{tab:ajuste}Índices de ajuste de los modelos ponderados}
\centering
\resizebox{\linewidth}{!}{\begin{tabular}[t]{lrrrlrrrrrrrrr}
\toprule
Modelo & n & Chi2 esc. & gl & p & CFI rob. & TLI rob. & RMSEA rob. & IC inf. & IC sup. & SRMR & CFI esc. & TLI esc. & RMSEA esc.\\
\midrule
AFC & 46.392 & 27.181,35 & 272 & <,001 & 0,900 & 0,890 & 0,092 & 0,091 & 0,093 & 0,040 & 0,981 & 0,979 & 0,046\\
SEM principal & 46.392 & 27.181,35 & 272 & <,001 & 0,900 & 0,890 & 0,092 & 0,091 & 0,093 & 0,040 & 0,981 & 0,979 & 0,046\\
Alternativo 2D & 46.392 & 23.275,26 & 269 & <,001 & 0,913 & 0,903 & 0,086 & 0,085 & 0,087 & 0,038 & 0,984 & 0,982 & 0,043\\
Rival inverso & 46.392 & 27.181,35 & 272 & <,001 & 0,900 & 0,890 & 0,092 & 0,091 & 0,093 & 0,040 & 0,981 & 0,979 & 0,046\\
Configural & 46.392 & 27.451,29 & 544 & <,001 & 0,898 & 0,888 & 0,093 & 0,091 & 0,094 & 0,041 & 0,981 & 0,979 & 0,046\\
\addlinespace
Umbrales & 46.392 & 27.778,41 & 582 & <,001 & NA & NA & NA & NA & NA & 0,041 & 0,981 & 0,980 & 0,045\\
Métrica & 46.392 & 27.091,71 & 604 & <,001 & NA & NA & NA & NA & NA & 0,041 & 0,981 & 0,981 & 0,043\\
\bottomrule
\end{tabular}}
\end{table}
\endgroup


El ajuste local mostró diez de 300 residuos de correlación con valor absoluto mayor que 0{,}10. El mayor residuo absoluto fue -0{,}219 entre los reactivos P2025S4 (desconocidos) y P2025S6 (amistades). Entre los índices de modificación más altos se encontraron las parejas P2025S5--P2025S6 y P5263S3--P5263S9, que podrían compartir contenido o contexto de respuesta. No se liberaron parámetros a partir de esos diagnósticos sin validación y razón sustantiva. Así se evita convertir una búsqueda exploratoria en una confirmación artificial del ajuste. La posible correlación residual P2025S5--P2025S6 queda como análisis futuro sujeto a calibración y validación en muestras separadas; no se estimó en este trabajo.

Las cargas estandarizadas del AFC oscilaron entre 0{,}651 y 0{,}753 para confianza interpersonal, 0{,}660 y 0{,}914 para confianza institucional, y 0{,}857 y 0{,}944 para familiaridad con mecanismos. Todas fueron positivas y estadísticamente distintas de cero bajo los errores estándar disponibles, aunque esos errores no incorporan el diseño complejo completo. Las cargas individuales se presentan en el apéndice reproducible.

\subsection{Confiabilidad y validez}

La confiabilidad compuesta fue 0{,}793 para confianza interpersonal, 0{,}968 para confianza institucional y 0{,}891 para familiaridad con mecanismos. Los valores superan 0{,}70, aunque la confiabilidad alta de la batería institucional debe leerse junto con su ajuste absoluto débil y con la heterogeneidad conceptual de las instituciones. Las AVE fueron 0{,}492, 0{,}668 y 0{,}794, respectivamente. La AVE interpersonal quedó apenas por debajo del referente 0{,}50; la evidencia convergente de ese factor es, por tanto, limítrofe.

La matriz de Fornell--Larcker muestra que las raíces de AVE (0{,}701, 0{,}817 y 0{,}891) superan las correlaciones latentes correspondientes (0{,}265, 0{,}111 y 0{,}090). Los HTMT no ponderados variaron entre 0{,}069 y 0{,}286, por debajo del criterio orientativo 0{,}85. Los criterios de discriminación son favorables en esos diagnósticos, pero no reparan el ajuste global débil. La Tabla~\ref{tab:validez} resume CR, AVE y máximos por factor; la Tabla~\ref{tab:discriminante} muestra ambas matrices completas. HTMT se calculó sin ponderar.
\begin{table}

\caption{\label{tab:discriminante}Matrices Fornell--Larcker y HTMT}
\centering
\begin{tabular}[t]{lrrrrrr}
\toprule
Constructo & FL-IP & FL-Inst. & FL-Fam. & HTMT-IP & HTMT-Inst. & HTMT-Fam.\\
\midrule
Interpersonal & 0,701 & 0,265 & 0,111 & 1,000 & 0,286 & 0,069\\
Institucional & 0,265 & 0,817 & 0,090 & 0,286 & 1,000 & 0,090\\
Familiaridad & 0,111 & 0,090 & 0,891 & 0,069 & 0,090 & 1,000\\
\bottomrule
\end{tabular}
\end{table}


\begin{table}

\caption{\label{tab:validez}Confiabilidad, AVE y validez discriminante}
\centering
\resizebox{\linewidth}{!}{\begin{tabular}[t]{lrrrrr}
\toprule
Constructo & CR & AVE & Raíz AVE & Cor. lat. máx. & HTMT máx. no pond.\\
\midrule
ConfianzaInterpersonal & 0,793 & 0,492 & 0,701 & 0,265 & 0,286\\
ConfianzaInstitucional & 0,968 & 0,668 & 0,817 & 0,265 & 0,286\\
ConocimientoMecanismos & 0,891 & 0,794 & 0,891 & 0,111 & 0,090\\
\bottomrule
\end{tabular}}
\end{table}


\subsection{Rutas estructurales e hipótesis}

La confianza interpersonal mostró una asociación positiva con la familiaridad (B=0{,}119, EE=0{,}011, \(z=11{,}05\), \(\beta=0{,}094\), IC 95\% para B [0{,}098,0{,}141], \(p<{,}001\), unilateral). La confianza institucional también se asoció positivamente (B=0{,}066, EE=0{,}008, \(z=8{,}19\), \(\beta=0{,}065\), IC 95\% [0{,}050,0{,}082], \(p<{,}001\), unilateral). En dirección y magnitud, ambos resultados son compatibles con H1 y H2; sin embargo, la baja magnitud, el ajuste débil y la varianza no ajustada al diseño impiden describirlos como confirmación poblacional definitiva. La covarianza estandarizada de los predictores fue 0{,}265 (H3), positiva y compatible con que las dimensiones estén relacionadas sin ser equivalentes.

El factor endógeno tuvo \(R^2=0{,}016\): el modelo explicó aproximadamente 1,6\% de la variación latente. Los tamaños incrementales únicos fueron \(f^2=.0083\) para confianza interpersonal y \(f^2=.0040\) para confianza institucional, ambos muy pequeños. Estos tamaños son coherentes con la recomendación de interpretar efectos y no confundir valores p pequeños, favorecidos por el N, con relevancia sustantiva.

\begin{table}

\caption{\label{tab:hipotesis}Hipótesis y parámetros estructurales}
\centering
\begin{tabular}[t]{llrrrrlrr}
\toprule
H & Ruta & B & EE & z & beta & p & IC\_inf & IC\_sup\\
\midrule
H1 & Interpersonal & 0,119 & 0,011 & 11,051 & 0,094 & <,001 & 0,098 & 0,141\\
H2 & Institucional & 0,066 & 0,008 & 8,190 & 0,065 & <,001 & 0,050 & 0,082\\
\bottomrule
\end{tabular}
\end{table}


El AFC y el SEM de referencia presentaron ajuste numéricamente equivalente (diferencia de grados de libertad igual a cero). El SEM no impone una restricción de covarianza latente que lo convierta en una prueba anidada más parsimoniosa: con dos predictores correlacionados y un resultado, sus parámetros estructurales reproducen una matriz latente libre. Por ello, la prueba de diferencia de chi-cuadrado no es interpretable como prueba de superioridad ni de deterioro estructural. La estructura alternativa de confianza institucional en dos factores tuvo CFI robusto=0{,}913, TLI=0{,}903, RMSEA robusto=0{,}086, IC 90\% [0{,}085,0{,}087] y SRMR=0{,}038. Aunque mejora algunos índices frente al modelo de tres factores, el ajuste robusto continúa débil; se mantiene como alternativa exploratoria, sin declarar un modelo ganador por ajuste únicamente. AIC y BIC no estuvieron disponibles con WLSMV.

El modelo rival invirtió la dirección, con familiaridad como predictor de confianza interpersonal e institucional, y dejó libre la covarianza residual entre estas últimas. Su ajuste global coincidió con el de referencia al redondear (\(\chi^2_s=27181{,}35\), 272 gl; CFI robusto=0{,}900; RMSEA robusto=0{,}092). En el rival, familiaridad se asoció con confianza interpersonal (\(B=0{,}088\), \(\beta=0{,}111\), EE=0{,}006) e institucional (\(B=0{,}088\), \(\beta=0{,}090\), EE=0{,}007), ambas con \(p<{,}001\). Esta equivalencia observacional impide elegir la dirección a partir del ajuste; la dirección del modelo principal proviene de la teoría y no de una superioridad empírica del ajuste \citep{MacCallum1993}.

No se calcularon efectos indirectos: la especificación no contiene mediador ni otro resultado endógeno. En consecuencia, los efectos directos son las dos rutas informadas, y el efecto total coincide con cada ruta directa en este modelo. No corresponde añadir intervalos bootstrap de mediación a una estructura que no la plantea.

\subsection{Sensibilidad a pesos y faltantes}
La comparación entre casos completos sin ponderación, casos completos con \texttt{FEX\_P} y la especificación principal por pares con pesos permite evaluar la sensibilidad de las rutas a estas decisiones. La Tabla~\ref{tab:sensibilidad} se genera desde las salidas del script estructural. La magnitud relativa de las asociaciones puede cambiar entre especificaciones; esto se presenta como sensibilidad metodológica y no como contradicción, dado que también cambia la muestra efectiva y el peso de cada observación.
\begin{table}

\caption{\label{tab:sensibilidad}Sensibilidad a pesos y tratamiento de faltantes}
\centering
\resizebox{\linewidth}{!}{\begin{tabular}[t]{llrrrlrrr}
\toprule
Especificacion & Ruta & B & EE & z & p & beta & IC\_inf & IC\_sup\\
\midrule
CC sin pesos & Interpersonal & 0,092 & 0,009 & 10,011 & <0,001 & 0,073 & 0,074 & 0,110\\
CC sin pesos & Institucional & 0,085 & 0,007 & 12,824 & <0,001 & 0,085 & 0,072 & 0,098\\
CC con FEX\_P & Interpersonal & 0,100 & 0,012 & 8,095 & <0,001 & 0,080 & 0,076 & 0,125\\
CC con FEX\_P & Institucional & 0,073 & 0,009 & 7,917 & <0,001 & 0,072 & 0,055 & 0,091\\
Pares con FEX\_P & Interpersonal & 0,119 & 0,011 & 11,051 & <0,001 & 0,094 & 0,098 & 0,141\\
\addlinespace
Pares con FEX\_P & Institucional & 0,066 & 0,008 & 8,190 & <0,001 & 0,065 & 0,050 & 0,082\\
\bottomrule
\end{tabular}}
% FUENTE_ECP_2023
\par\smallskip\noindent\parbox{\linewidth}{\footnotesize\textit{Fuente:} elaboración propia con microdatos de la ECP 2023 del DANE.}
\end{table}


\subsection{Sensibilidad aproximada al efecto de diseño}
Como diagnóstico conservador, se aproximó el efecto de diseño mediante Kish, $deff=1+CV(FEX\_P)^2$, y se inflaron los errores estándar por $\sqrt{deff}$. Este cálculo solo refleja la desigualdad de pesos; no modela conglomeración, estratificación ni faltantes y no sustituye una corrección completa del diseño. Las conclusiones sobre significación deben matizarse con este análisis exploratorio.
El efecto aproximado fue \(deff=1{,}758\). Tras inflar los EE, las rutas mantuvieron sus signos y su evidencia estadística (interpersonal: \(EE=0{,}014\), \(z=8{,}33\); institucional: \(EE=0{,}011\), \(z=6{,}18\); ambas \(p<{,}001\)). Esto no valida los errores estándar del diseño, pues la aproximación no incorpora sus componentes no disponibles.
\begin{table}

\caption{\label{tab:deff}Sensibilidad aproximada al efecto de diseño por desigualdad de pesos}
\centering
\resizebox{\linewidth}{!}{\begin{tabular}[t]{lrrrrrlrr}
\toprule
Ruta & B & EE\_original & deff & EE\_inflado & z & p & IC\_inf & IC\_sup\\
\midrule
Interpersonal & 0,119 & 0,011 & 1,758 & 0,014 & 8,334 & <,001 & 0,091 & 0,147\\
Institucional & 0,066 & 0,008 & 1,758 & 0,011 & 6,177 & <,001 & 0,045 & 0,087\\
\bottomrule
\end{tabular}}
\end{table}


\subsection{Invariancia y diferencias de rutas por sexo}

Los dos grupos superaron ampliamente el tamaño mínimo operativo de 200: 21.405 hombres y 24.987 mujeres. En la secuencia ordinal, el CFI escalado fue 0{,}981 en el modelo configural, 0{,}981 con umbrales y 0{,}981 con umbrales y cargas; los cambios fueron -0{,}0002 y +0{,}0005. El RMSEA escalado pasó de 0{,}046 a 0{,}045 y 0{,}044, y el SRMR permaneció cercano a 0{,}041. Sin embargo, la prueba de diferencia configural--umbrales fue significativa (\(\Delta\chi^2=59{,}51\), 38 gl, \(p=0{,}014\)); la de umbrales--métrica no fue significativa (\(\Delta\chi^2=31{,}32\), 22 gl, \(p=0{,}090\)). Los cambios aproximados son pequeños, pero la discrepancia de chi-cuadrado no permite presentar la secuencia como evidencia inequívoca. Además, el ajuste robusto absoluto del modelo configural fue débil (CFI=0{,}898; RMSEA=0{,}093).

Al imponer conjuntamente igualdad de rutas, la prueba ajustada fue significativa (\(\Delta\chi^2=8{,}21\), 2 gl, \(p=0{,}017\)). La prueba Wald por ruta no indicó diferencia en la asociación interpersonal (\(\beta_{H}=0{,}096\), \(\beta_{M}=0{,}090\), \(p=0{,}650\)). La asociación institucional fue mayor en hombres (\(\beta=0{,}083\)) que en mujeres (\(\beta=0{,}049\)); la prueba de igualdad dio \(p=0{,}034\). Al ajustar las dos pruebas de Wald, H2 no conserva significación con Bonferroni (umbral 0{,}025) ni con Holm; la prueba conjunta de rutas (\(p=0{,}017\)) es un contraste distinto y sí permanece significativa sin este ajuste de dos pruebas. En consecuencia, la diferencia institucional por sexo es un hallazgo exploratorio y matizado. Como las comparaciones de rutas se basan en equivalencia métrica aproximada y errores que no incorporan UPM y estratos, se interpretan como indicios de posible moderación asociativa.

\begin{table}

\caption{\label{tab:invarianza}Pruebas de diferencia y cambios de ajuste entre niveles de invariancia}
\centering
\resizebox{\linewidth}{!}{\begin{tabular}[t]{lllllll}
\toprule
Restricción añadida & Delta Chi2 & Delta gl & p & Delta CFI & Delta RMSEA & Delta SRMR\\
\midrule
Umbrales & 59,51 & 38 & 0,014 & -0,0002 & -0,0013 & 0,0000\\
Cargas y umbrales & 31,32 & 22 & 0,090 & 0,0005 & -0,0014 & 0,0000\\
\bottomrule
\end{tabular}}
% FUENTE_ECP_2023
\par\smallskip\noindent\parbox{\linewidth}{\footnotesize\textit{Fuente:} elaboración propia con microdatos de la ECP 2023 del DANE.}
\end{table}

La prueba configural--umbrales fue significativa (\(p=0{,}014\)) y el paso umbrales--métrica no lo fue (\(p=0{,}090\)). Los cambios de ajuste escalado fueron pequeños, pero no eliminan la discrepancia con las pruebas de diferencia. Los índices robustos de algunos modelos multigrupo no están definidos por lavaan y se reportan como no disponibles.

\begin{table}

\caption{\label{tab:waldsexo}Pruebas de Wald por ruta con ajuste por comparaciones múltiples}
\centering
\begin{tabular}[t]{lrrrrr}
\toprule
Hipotesis & Wald & gl & p & p\_Bonferroni & p\_Holm\\
\midrule
H1 & 0,206 & 1 & 0,650 & 1,000 & 0,650\\
H2 & 4,512 & 1 & 0,034 & 0,067 & 0,067\\
\bottomrule
\end{tabular}
\end{table}


\section{Discusión}

Este estudio examinó la relación entre confianza y familiaridad con mecanismos formales de participación usando la intersección de cinco módulos de la ECP 2023. Las dos rutas estimadas fueron positivas, como se planteó en H1 y H2, pero su magnitud fue pequeña y el modelo explicó una fracción reducida de la varianza. La significación de las pruebas no debe convertirse en una afirmación de importancia práctica ni en evidencia causal. En una muestra de este tamaño, incluso asociaciones modestas pueden producir valores p reducidos.

El resultado para confianza interpersonal es compatible con la idea de que los vínculos y la circulación de información se relacionan con reconocimiento cívico. No obstante, la batería agrupa confianza hacia personas desconocidas, vecinos, amistades y personas de otra nacionalidad. La asociación no permite concluir que la confianza generalizada por sí sola esté vinculada con la familiaridad. Sería útil distinguir en estudios posteriores confianza generalizada y particularizada, ampliar la medición y evaluar validez en una muestra independiente.

La asociación positiva de confianza institucional con familiaridad es coherente con una mayor orientación hacia canales institucionales, pero también admite explicaciones alternativas. Las personas más interesadas en política pueden conocer instituciones y mecanismos y, al mismo tiempo, formar evaluaciones particulares de ellas; variables como educación, exposición a trámites, edad, uso de medios o experiencia de participación podrían explicar parte de la asociación. El modelo no incluye esos factores ni identifica su orden temporal. Por tanto, la ruta no significa que aumentar confianza vaya a producir mayor familiaridad.

La diferencia por sexo de la asociación institucional requiere especial prudencia. La magnitud estandarizada fue mayor entre hombres, pero la prueba por ruta es exploratoria y los errores no están ajustados a las unidades de muestreo. Además, la comparación supone invariancia de medición ordinal que muestra evidencia mixta: los cambios de ajuste aproximados son pequeños, mientras la prueba de umbrales rechaza igualdad al nivel convencional. La presencia de no aplicabilidad estructural en dos reactivos institucionales en Bogotá y el ajuste robusto configural débil también pueden influir en la comparabilidad. Un análisis de seguimiento debería examinar invariancia parcial con criterios sustantivos, datos adicionales de diseño y replicación independiente antes de sostener una diferencia poblacional.

El principal resultado metodológico es que la medición no puede darse por resuelta. El ajuste robusto del AFC no alcanza los criterios orientativos usuales de CFI/TLI y RMSEA, aunque sus versiones escaladas son más favorables y el SRMR es bajo. La discrepancia entre familias de índices para WLSMV ha sido documentada y no se resuelve seleccionando solo la cifra más favorable \citep{XiaYang2019}. La batería institucional puede estar representando múltiples dimensiones de evaluación pública, y los residuos/IM detectan dependencias locales que conviene explicar mediante contenido. La alternativa de dos factores mejora parcialmente el ajuste robusto, pero sigue siendo exploratoria y necesita respaldo conceptual. No se añadieron correlaciones residuales post hoc porque los índices de modificación, por sí solos, no justifican cambios confirmatorios.

En confiabilidad, los valores compuestos fueron adecuados, pero la AVE interpersonal quedó marginalmente bajo 0{,}50. La evaluación discriminante resultó favorable en las correlaciones latentes y HTMT calculado, aunque este último no incorporó la ponderación. Por tanto, no debe afirmarse que la validez discriminante ponderada haya quedado demostrada. La conjunción de buen CR y ajuste débil muestra que consistencia interna alta no equivale a estructura factorial correcta.

El factor de expansión es necesario para ponderar los puntos estimados, y su incorporación diferencia este análisis de las estimaciones previas sin peso. No obstante, ponderar no reproduce automáticamente el diseño multietápico. Los identificadores anonimizados disponibles no permiten agrupar observaciones en las unidades de muestreo ni especificar estratos. Como el DANE no proporciona esos identificadores o varianzas alternativas para este análisis, los intervalos y pruebas no incorporan toda la incertidumbre muestral. La suma de los factores en la intersección tampoco se interpreta como un total poblacional oficial. La inferencia se limita a resultados ponderados exploratorios, y cualquier generalización formal debe reconocer esta restricción.

La definición del desenlace es otra delimitación importante. Se midió si la persona declara conocer o haber oído de mecanismos, no su conocimiento factual ni su uso. Las seis respuestas binarias son indicadores de reconocimiento y pueden reflejar niveles de exposición distintos. Una siguiente investigación podría contrastar este factor con preguntas de conocimiento factual, participación efectiva, educación cívica y experiencia de contacto institucional. También podría incorporar covariables justificadas y separar la medición por zona o grupo de edad si los tamaños y el diseño lo permiten.

Finalmente, los modelos de referencia e inverso pueden ser observacionalmente equivalentes cuando la matriz de covarianzas latentes es libre. La encuesta transversal no aporta la temporalidad o intervención necesarias para resolver la dirección. La discusión, por ello, se restringe a compatibilidad de asociaciones con las hipótesis, no a efectos ni mecanismos causales.

\subsection{Limitaciones}

Primero, el corte transversal impide establecer precedencia temporal o causalidad. Segundo, el ajuste robusto del modelo de tres factores fue débil y la AVE de confianza interpersonal resultó limítrofe; los índices escalados más favorables no eliminan esta preocupación. Tercero, el diseño complejo no pudo incorporarse plenamente a la varianza: se usó \texttt{FEX\_P}, pero faltan estratos y UPM/conglomerados y no se dispuso de varianzas calculadas por el DANE. Cuarto, los faltantes se trataron por pares y existen dos reactivos con no aplicabilidad territorial; la base analítica es una intersección de módulos y no el total oficial de encuestados. Quinto, HTMT fue calculado sin ponderación y funciona como diagnóstico complementario. Sexto, no se ajustaron covariables potencialmente asociadas con confianza y familiaridad. Finalmente, las medidas son autoinformadas y pueden compartir sesgo de método común.

\section{Conclusiones}

En la base integrada de la ECP 2023, la confianza interpersonal y la confianza institucional mostraron asociaciones positivas, pequeñas, con la familiaridad autorreportada con mecanismos de participación. El modelo explicó 1,6\% de la varianza del factor endógeno. La ruta institucional fue más alta entre hombres que entre mujeres en las pruebas multigrupo, mientras que no se detectó diferencia por sexo para la ruta interpersonal. Estos resultados responden la pregunta en términos asociativos, pero no prueban causalidad ni relevancia práctica elevada.

La conclusión sustantiva debe condicionarse al ajuste robusto débil y a la limitación de diseño: los pesos \texttt{FEX\_P} se incorporaron, pero no se pudieron corregir errores e intervalos por estratificación y conglomeración. Por esto, los contrastes y las diferencias de grupo son exploratorios. El estudio aporta una primera estimación reproducible y señala la necesidad de mejorar la medición, contrastar los hallazgos en otra muestra y disponer de componentes de varianza para inferencia plenamente ajustada.

\section*{Disponibilidad de datos y código}

Los microdatos anonimizados de la ECP 2023 se descargan del catálogo DANE-ANDA: \url{https://microdatos.dane.gov.co/index.php/catalog/822}. Los datos individuales no se redistribuyen en este paquete; el archivo \texttt{README.md} explica la descarga y organización requeridas. El código, resultados agregados y materiales LaTeX de este análisis están disponibles en el repositorio del proyecto: \url{https://github.com/juandavidroa06/SEM}. La tabla reproducible y los archivos de sesión se incluyen en el paquete de entrega.

\bibliographystyle{apalike}
\bibliography{referencias}

\appendix
\section{Estadísticos descriptivos por indicador}
\label{ap:descriptivos}
Los descriptivos corresponden a la base analítica integrada, sin ponderación, antes de la estimación WLSMV. En los reactivos dicotómicos, la media es la proporción no ponderada que reportó conocer u oír hablar del mecanismo.

\begin{longtable}[t]{lrrrrr}
\caption{Estadísticos descriptivos por indicador}\\
\toprule
Indicador & n válido & Faltantes \% & Media & DE & Mediana\\
\midrule
\endfirsthead
\toprule
Indicador & n válido & Faltantes \% & Media & DE & Mediana\\
\midrule
\endhead
P2025S4 & 45.690 & 1,513 & 1,716 & 0,958 & 1\\
P2025S5 & 45.911 & 1,037 & 2,808 & 1,030 & 3\\
P2025S6 & 45.967 & 0,916 & 3,265 & 1,010 & 3\\
P2025S7 & 45.288 & 2,380 & 1,767 & 0,973 & 1\\
P5263S1 & 41.447 & 10,659 & 2,566 & 1,062 & 3\\
\addlinespace
P5263S2 & 43.301 & 6,663 & 2,502 & 1,055 & 3\\
P5263S3 & 43.519 & 6,193 & 2,880 & 1,109 & 3\\
P5263S4 & 42.833 & 7,672 & 2,311 & 1,054 & 2\\
P5263S5 & 42.638 & 8,092 & 2,504 & 1,035 & 3\\
P5263S6 & 37.143 & 19,937 & 2,431 & 0,999 & 3\\
\addlinespace
P5263S7 & 41.624 & 10,278 & 2,553 & 1,034 & 3\\
P5263S8 & 42.399 & 8,607 & 2,434 & 1,002 & 3\\
P5263S9 & 43.878 & 5,419 & 2,663 & 1,111 & 3\\
P5263S10 & 43.466 & 6,307 & 2,142 & 0,998 & 2\\
P5263S11 & 42.655 & 8,055 & 2,823 & 1,062 & 3\\
\addlinespace
P5263S12 & 38.203 & 17,652 & 2,500 & 1,035 & 3\\
P5263S13 & 42.803 & 7,736 & 2,640 & 1,076 & 3\\
P5263S14 & 43.100 & 7,096 & 2,836 & 1,090 & 3\\
P5263S15 & 43.564 & 6,096 & 2,511 & 1,111 & 3\\
P5376S1 & 46.392 & 0,000 & 0,411 & 0,492 & 0\\
\addlinespace
P5376S2 & 46.392 & 0,000 & 0,230 & 0,421 & 0\\
P5376S3 & 46.392 & 0,000 & 0,398 & 0,490 & 0\\
P5376S4 & 46.392 & 0,000 & 0,426 & 0,494 & 0\\
P5376S5 & 46.392 & 0,000 & 0,378 & 0,485 & 0\\
P5376S6 & 46.392 & 0,000 & 0,250 & 0,433 & 0\\
\bottomrule
\end{longtable}
% FUENTE_ECP_2023
\par\smallskip\noindent\parbox{\linewidth}{\footnotesize\textit{Fuente:} elaboración propia con microdatos de la ECP 2023 del DANE.}


\section{Cargas estandarizadas del modelo de medición}
Las cargas, errores estándar y estadísticos $z$ del AFC ponderado se presentan para facilitar la evaluación de cada indicador.

\begin{longtable}[t]{llrrrlrr}
\caption{Cargas factoriales estandarizadas del AFC ponderado}\\
\toprule
Factor & Indicador & Carga est. & EE est. & z & p & IC inf. & IC sup.\\
\midrule
\endfirsthead
\toprule
Factor & Indicador & Carga est. & EE est. & z & p & IC inf. & IC sup.\\
\midrule
\endhead
Interpersonal & P2025S4 & 0,674 & 0,007 & 99,960 & <0,001 & 0,661 & 0,688\\
Interpersonal & P2025S5 & 0,753 & 0,006 & 133,547 & <0,001 & 0,742 & 0,764\\
Interpersonal & P2025S6 & 0,652 & 0,006 & 113,211 & <0,001 & 0,640 & 0,663\\
Interpersonal & P2025S7 & 0,722 & 0,007 & 108,813 & <0,001 & 0,709 & 0,736\\
Institucional & P5263S1 & 0,840 & 0,002 & 461,004 & <0,001 & 0,837 & 0,844\\
\addlinespace
Institucional & P5263S2 & 0,790 & 0,002 & 352,675 & <0,001 & 0,786 & 0,794\\
Institucional & P5263S3 & 0,720 & 0,003 & 251,900 & <0,001 & 0,714 & 0,726\\
Institucional & P5263S4 & 0,841 & 0,002 & 462,356 & <0,001 & 0,837 & 0,844\\
Institucional & P5263S5 & 0,869 & 0,002 & 575,806 & <0,001 & 0,866 & 0,872\\
Institucional & P5263S6 & 0,914 & 0,001 & 713,081 & <0,001 & 0,912 & 0,917\\
\addlinespace
Institucional & P5263S7 & 0,899 & 0,001 & 720,036 & <0,001 & 0,897 & 0,901\\
Institucional & P5263S8 & 0,888 & 0,001 & 647,857 & <0,001 & 0,885 & 0,890\\
Institucional & P5263S9 & 0,741 & 0,003 & 277,267 & <0,001 & 0,735 & 0,746\\
Institucional & P5263S10 & 0,773 & 0,002 & 314,522 & <0,001 & 0,768 & 0,777\\
Institucional & P5263S11 & 0,806 & 0,002 & 387,016 & <0,001 & 0,801 & 0,810\\
\addlinespace
Institucional & P5263S12 & 0,862 & 0,002 & 462,653 & <0,001 & 0,859 & 0,866\\
Institucional & P5263S13 & 0,856 & 0,002 & 523,336 & <0,001 & 0,853 & 0,859\\
Institucional & P5263S14 & 0,758 & 0,003 & 301,571 & <0,001 & 0,753 & 0,763\\
Institucional & P5263S15 & 0,660 & 0,003 & 204,364 & <0,001 & 0,654 & 0,667\\
Familiaridad & P5376S1 & 0,857 & 0,003 & 251,227 & <0,001 & 0,850 & 0,863\\
\addlinespace
Familiaridad & P5376S2 & 0,944 & 0,003 & 321,021 & <0,001 & 0,938 & 0,950\\
Familiaridad & P5376S3 & 0,899 & 0,003 & 326,149 & <0,001 & 0,894 & 0,904\\
Familiaridad & P5376S4 & 0,893 & 0,003 & 308,940 & <0,001 & 0,887 & 0,899\\
Familiaridad & P5376S5 & 0,874 & 0,003 & 278,367 & <0,001 & 0,868 & 0,881\\
Familiaridad & P5376S6 & 0,876 & 0,004 & 227,198 & <0,001 & 0,869 & 0,884\\
\bottomrule
\end{longtable}
% FUENTE_ECP_2023
\par\smallskip\noindent\parbox{\linewidth}{\footnotesize\textit{Fuente:} elaboración propia con microdatos de la ECP 2023 del DANE.}


\section{Matriz policórica de los 25 indicadores}
\label{ap:policoricas}
La matriz fue calculada en la Fase 1 con los 35.024 casos completos y códigos ordinales/binarios, sin ponderación; se adjunta para trazabilidad descriptiva y no como matriz ponderada del modelo WLSMV de Fase 2. Las etiquetas son los códigos del formulario.
\begin{landscape}

\begingroup\tiny\setlength{\tabcolsep}{1.2pt}\begin{longtable}[t]{lrrrrrrrrrrrrrrrrrrrrrrrrr}
\caption{Matriz policórica no ponderada}\\
\toprule
Indicador & I4 & I5 & I6 & I7 & T1 & T2 & T3 & T4 & T5 & T6 & T7 & T8 & T9 & T10 & T11 & T12 & T13 & T14 & T15 & K1 & K2 & K3 & K4 & K5 & K6\\
\midrule
\endfirsthead
\toprule
Indicador & I4 & I5 & I6 & I7 & T1 & T2 & T3 & T4 & T5 & T6 & T7 & T8 & T9 & T10 & T11 & T12 & T13 & T14 & T15 & K1 & K2 & K3 & K4 & K5 & K6\\
\midrule
\endhead
I4 & 1,00 & 0,42 & 0,26 & 0,60 & 0,13 & 0,18 & 0,08 & 0,23 & 0,17 & 0,18 & 0,14 & 0,18 & 0,14 & 0,24 & 0,09 & 0,17 & 0,13 & 0,07 & 0,16 & 0,03 & 0,08 & 0,03 & 0,00 & 0,01 & 0,04\\
I5 & 0,42 & 1,00 & 0,64 & 0,40 & 0,14 & 0,16 & 0,13 & 0,18 & 0,16 & 0,16 & 0,15 & 0,16 & 0,16 & 0,19 & 0,16 & 0,16 & 0,14 & 0,12 & 0,17 & 0,06 & 0,06 & 0,03 & 0,00 & -0,01 & 0,09\\
I6 & 0,26 & 0,64 & 1,00 & 0,35 & 0,13 & 0,15 & 0,12 & 0,14 & 0,16 & 0,15 & 0,15 & 0,15 & 0,14 & 0,16 & 0,16 & 0,15 & 0,13 & 0,12 & 0,16 & 0,12 & 0,08 & 0,09 & 0,08 & 0,09 & 0,13\\
I7 & 0,60 & 0,40 & 0,35 & 1,00 & 0,15 & 0,20 & 0,11 & 0,22 & 0,18 & 0,20 & 0,17 & 0,20 & 0,14 & 0,24 & 0,15 & 0,19 & 0,16 & 0,12 & 0,18 & 0,10 & 0,10 & 0,08 & 0,05 & 0,05 & 0,09\\
T1 & 0,13 & 0,14 & 0,13 & 0,15 & 1,00 & 0,72 & 0,66 & 0,71 & 0,75 & 0,74 & 0,82 & 0,71 & 0,61 & 0,62 & 0,69 & 0,71 & 0,76 & 0,67 & 0,53 & 0,10 & 0,10 & 0,12 & 0,13 & 0,12 & 0,13\\
\addlinespace
T2 & 0,18 & 0,16 & 0,15 & 0,20 & 0,72 & 1,00 & 0,59 & 0,75 & 0,71 & 0,76 & 0,72 & 0,79 & 0,60 & 0,68 & 0,63 & 0,74 & 0,67 & 0,59 & 0,59 & 0,05 & 0,07 & 0,07 & 0,06 & 0,04 & 0,08\\
T3 & 0,08 & 0,13 & 0,12 & 0,11 & 0,66 & 0,59 & 1,00 & 0,57 & 0,63 & 0,62 & 0,66 & 0,61 & 0,70 & 0,49 & 0,64 & 0,62 & 0,64 & 0,63 & 0,44 & 0,11 & 0,09 & 0,12 & 0,11 & 0,13 & 0,12\\
T4 & 0,23 & 0,18 & 0,14 & 0,22 & 0,71 & 0,75 & 0,57 & 1,00 & 0,79 & 0,81 & 0,75 & 0,79 & 0,62 & 0,76 & 0,61 & 0,73 & 0,70 & 0,57 & 0,62 & 0,01 & 0,06 & 0,03 & 0,03 & 0,03 & 0,08\\
T5 & 0,17 & 0,16 & 0,16 & 0,18 & 0,75 & 0,71 & 0,63 & 0,79 & 1,00 & 0,82 & 0,83 & 0,77 & 0,65 & 0,69 & 0,70 & 0,74 & 0,78 & 0,66 & 0,61 & 0,07 & 0,09 & 0,09 & 0,10 & 0,09 & 0,10\\
T6 & 0,18 & 0,16 & 0,15 & 0,20 & 0,74 & 0,76 & 0,62 & 0,81 & 0,82 & 1,00 & 0,83 & 0,87 & 0,63 & 0,75 & 0,69 & 0,81 & 0,74 & 0,65 & 0,63 & 0,04 & 0,07 & 0,07 & 0,08 & 0,06 & 0,09\\
\addlinespace
T7 & 0,14 & 0,15 & 0,15 & 0,17 & 0,82 & 0,72 & 0,66 & 0,75 & 0,83 & 0,83 & 1,00 & 0,81 & 0,66 & 0,68 & 0,75 & 0,78 & 0,81 & 0,71 & 0,58 & 0,09 & 0,10 & 0,11 & 0,11 & 0,11 & 0,11\\
T8 & 0,18 & 0,16 & 0,15 & 0,20 & 0,71 & 0,79 & 0,61 & 0,79 & 0,77 & 0,87 & 0,81 & 1,00 & 0,64 & 0,75 & 0,68 & 0,81 & 0,72 & 0,64 & 0,62 & 0,05 & 0,07 & 0,07 & 0,07 & 0,06 & 0,09\\
T9 & 0,14 & 0,16 & 0,14 & 0,14 & 0,61 & 0,60 & 0,70 & 0,62 & 0,65 & 0,63 & 0,66 & 0,64 & 1,00 & 0,60 & 0,62 & 0,63 & 0,66 & 0,58 & 0,48 & 0,06 & 0,06 & 0,07 & 0,07 & 0,08 & 0,08\\
T10 & 0,24 & 0,19 & 0,16 & 0,24 & 0,62 & 0,68 & 0,49 & 0,76 & 0,69 & 0,75 & 0,68 & 0,75 & 0,60 & 1,00 & 0,57 & 0,72 & 0,64 & 0,51 & 0,61 & 0,02 & 0,07 & 0,05 & 0,05 & 0,04 & 0,10\\
T11 & 0,09 & 0,16 & 0,16 & 0,15 & 0,69 & 0,63 & 0,64 & 0,61 & 0,70 & 0,69 & 0,75 & 0,68 & 0,62 & 0,57 & 1,00 & 0,72 & 0,77 & 0,75 & 0,57 & 0,08 & 0,07 & 0,11 & 0,10 & 0,09 & 0,09\\
\addlinespace
T12 & 0,17 & 0,16 & 0,15 & 0,19 & 0,71 & 0,74 & 0,62 & 0,73 & 0,74 & 0,81 & 0,78 & 0,81 & 0,63 & 0,72 & 0,72 & 1,00 & 0,77 & 0,69 & 0,63 & 0,06 & 0,08 & 0,09 & 0,08 & 0,07 & 0,09\\
T13 & 0,13 & 0,14 & 0,13 & 0,16 & 0,76 & 0,67 & 0,64 & 0,70 & 0,78 & 0,74 & 0,81 & 0,72 & 0,66 & 0,64 & 0,77 & 0,77 & 1,00 & 0,75 & 0,57 & 0,07 & 0,09 & 0,10 & 0,09 & 0,09 & 0,08\\
T14 & 0,07 & 0,12 & 0,12 & 0,12 & 0,67 & 0,59 & 0,63 & 0,57 & 0,66 & 0,65 & 0,71 & 0,64 & 0,58 & 0,51 & 0,75 & 0,69 & 0,75 & 1,00 & 0,52 & 0,08 & 0,05 & 0,10 & 0,09 & 0,09 & 0,09\\
T15 & 0,16 & 0,17 & 0,16 & 0,18 & 0,53 & 0,59 & 0,44 & 0,62 & 0,61 & 0,63 & 0,58 & 0,62 & 0,48 & 0,61 & 0,57 & 0,63 & 0,57 & 0,52 & 1,00 & 0,00 & 0,04 & 0,03 & 0,04 & 0,03 & 0,07\\
K1 & 0,03 & 0,06 & 0,12 & 0,10 & 0,10 & 0,05 & 0,11 & 0,01 & 0,07 & 0,04 & 0,09 & 0,05 & 0,06 & 0,02 & 0,08 & 0,06 & 0,07 & 0,08 & 0,00 & 1,00 & 0,84 & 0,80 & 0,72 & 0,70 & 0,72\\
\addlinespace
K2 & 0,08 & 0,06 & 0,08 & 0,10 & 0,10 & 0,07 & 0,09 & 0,06 & 0,09 & 0,07 & 0,10 & 0,07 & 0,06 & 0,07 & 0,07 & 0,08 & 0,09 & 0,05 & 0,04 & 0,84 & 1,00 & 0,85 & 0,79 & 0,85 & 0,82\\
K3 & 0,03 & 0,03 & 0,09 & 0,08 & 0,12 & 0,07 & 0,12 & 0,03 & 0,09 & 0,07 & 0,11 & 0,07 & 0,07 & 0,05 & 0,11 & 0,09 & 0,10 & 0,10 & 0,03 & 0,80 & 0,85 & 1,00 & 0,85 & 0,77 & 0,74\\
K4 & 0,00 & 0,00 & 0,08 & 0,05 & 0,13 & 0,06 & 0,11 & 0,03 & 0,10 & 0,08 & 0,11 & 0,07 & 0,07 & 0,05 & 0,10 & 0,08 & 0,09 & 0,09 & 0,04 & 0,72 & 0,79 & 0,85 & 1,00 & 0,79 & 0,78\\
K5 & 0,01 & -0,01 & 0,09 & 0,05 & 0,12 & 0,04 & 0,13 & 0,03 & 0,09 & 0,06 & 0,11 & 0,06 & 0,08 & 0,04 & 0,09 & 0,07 & 0,09 & 0,09 & 0,03 & 0,70 & 0,85 & 0,77 & 0,79 & 1,00 & 0,79\\
K6 & 0,04 & 0,09 & 0,13 & 0,09 & 0,13 & 0,08 & 0,12 & 0,08 & 0,10 & 0,09 & 0,11 & 0,09 & 0,08 & 0,10 & 0,09 & 0,09 & 0,08 & 0,09 & 0,07 & 0,72 & 0,82 & 0,74 & 0,78 & 0,79 & 1,00\\
\bottomrule
\end{longtable}
% FUENTE_ECP_2023
\par\smallskip\noindent\parbox{\linewidth}{\footnotesize\textit{Fuente:} elaboración propia con microdatos de la ECP 2023 del DANE.}
\endgroup

\end{landscape}

\section{Sintaxis reproducible del análisis}
\label{ap:codigo}
Los scripts de integración, estimación, diagnóstico e invariancia se ejecutan en orden de \texttt{00\_preparacion.R} a \texttt{04\_exportar\_resultados.R}. Requieren los cinco ZIP de módulos y \texttt{Tabla de viviendas.zip} en la carpeta relativa \texttt{datos\_crudos/}, descargados por separado del catálogo DANE. La sintaxis del modelo de referencia es:
\begin{verbatim}
modelo_medicion <- '
  ConfianzaInterpersonal =~ P2025S4 + P2025S5 + P2025S6 + P2025S7
  ConfianzaInstitucional =~ P5263S1 + P5263S2 + P5263S3 + P5263S4 + P5263S5 +
                            P5263S6 + P5263S7 + P5263S8 + P5263S9 + P5263S10 +
                            P5263S11 + P5263S12 + P5263S13 + P5263S14 + P5263S15
  ConocimientoMecanismos =~ P5376S1 + P5376S2 + P5376S3 + P5376S4 + P5376S5 + P5376S6
  ConocimientoMecanismos ~ ConfianzaInterpersonal + ConfianzaInstitucional
  ConfianzaInterpersonal ~~ ConfianzaInstitucional
'
fit <- lavaan::sem(modelo_sem, data=datos, ordered=items,
                   estimator='WLSMV', missing='pairwise',
                   sampling.weights='FEX_P')
\end{verbatim}

\subsection*{Scripts reproducibles}
La secuencia de ejecución es \texttt{00\_preparacion.R}, \texttt{01\_afc.R}, \texttt{02\_estructural.R}, \texttt{03\_multigrupo.R} y \texttt{04\_exportar\_resultados.R}. Los archivos completos y las salidas agregadas acompañan el manuscrito.

\end{document}

